% Generated draft: requires mathematical and source review.
\documentclass[11pt]{article}
\usepackage[margin=25mm]{geometry}
\usepackage{fontspec}
\setmainfont{Latin Modern Roman}
\usepackage{amsmath,amssymb,mathtools,unicode-math}
\setmathfont{Latin Modern Math}
\usepackage{enumitem,needspace,xurl,hyperref}
\setlength{\parindent}{0pt}
\setlength{\parskip}{5pt}
\setlength{\emergencystretch}{3em}
\newcommand{\N}{\mathbb{N}}
\newcommand{\Z}{\mathbb{Z}}
\newcommand{\Q}{\mathbb{Q}}
\newcommand{\R}{\mathbb{R}}
\newcommand{\C}{\mathbb{C}}
\newcommand{\F}{\mathbb{F}}
\newcommand{\A}{\mathbb{A}}
\newcommand{\PP}{\mathbb{P}}
\newcommand{\E}{\mathbb{E}}
\newcommand{\eps}{\varepsilon}
\newcommand{\dd}{\,\mathrm{d}}
\begin{document}
\section*{UnsolvedMath ID: 30000748}
\textbf{Source problem code:} catalogue entry
\subsection{Definitions and mathematical statement}
This record asks two separate positional-game questions.

First, does a Breaker winning strategy in the intended original Maker--Breaker game imply a winning strategy in its recycled version, and does this hold in particular when the original win is guaranteed by the Erd\H{o}s--Selfridge theorem? In an ordinary hypergraph Maker--Breaker game, Maker seeks to occupy a whole winning set and Breaker seeks to prevent this. The recycled rules, board, winning family, allowed biases and starting-player conventions are not supplied. The wording begins ``in the recycled game'' and then asks about implication to that same game, leaving the original-versus-recycled comparison implicit. It also does not explicitly identify the beneficiary of the second winning strategy, although Breaker is the natural reading. The precise implication must therefore be recovered rather than treated as a tautology or supplied with an invented recycling operation. The Erd\H{o}s--Selfridge subcase is retained without assuming a missing version of that theorem or its game hypotheses.

Second, consider the stated Kaplansky four-in-a-line game on \(\Z^2\). Under its one-point-per-turn reading, two players alternately claim an unclaimed lattice point. At a finite stage let \(S_1,S_2\subset\Z^2\) be their disjoint claimed sets. Player \(i\) wins upon creating an affine line \(\ell\subset\R^2\) such that
\[
 |S_i\cap\ell|\geq4,\qquad S_{3-i}\cap\ell=\varnothing.
\]
The line is opponent-free in its entirety, not just on the segment between the four points. The points are not required to be consecutive or equally spaced, and no horizontal/vertical/diagonal-only direction restriction is added. The first player to meet the condition wins.

The question is whether this infinite-board game is a draw under optimal play. This does not mean every possible play is drawn, and there is no requirement to fill the board in order to draw. A strategy must respond to all legal opposing play; the relevant draw claim concerns the absence of a forced finite win, with the intended defensive-strategy convention. The source does not spell out conventions for infinite play or passing, so those remain to be confirmed.

The literature triage checked 2026-08-21 says the Kaplansky game is still recorded as an unsolved draw question. It does not separately assess the recycled implication. The open label and reported source verification remain dated metadata, not independently checked game-theoretic results.
\subsection{Short English statement}
Ask whether Breaker strategies transfer to the undefined recycled game, including the Erd\H{o}s--Selfridge subcase, and separately whether Kaplansky's infinite four-in-a-line game is drawn. The latter uses any affine line containing four owned points and no opponent point, without a consecutiveness restriction. The recycled rules are missing.
\subsection{Sources}
\begin{itemize}
\item Local UnsolvedMath registry, permanent ID 30000748, source code catalogue entry, ``Recycled Maker--Breaker Games and Kaplansky Four-in-a-Line,'' including both original questions and the explicit opponent-free-line condition.
\item \emph{Mini-Workshop: Positional Games (2008)}, Oberwolfach Reports, pp.\ 1073--1114, DOI 10.4171/owr/2007/20: \url{https://doi.org/10.4171/owr/2007/20}. The citation uses 2008, while the DOI and proposed-year metadata use 2007.
\item Supplied reference ``Beck combinatorial games'': \url{https://www.nzdr.ru/data/media/biblio/kolxoz/M/MA/MAc/Beck%20J.%20Combinatorial%20games%20%28CUP%2C%202008%29%28ISBN%200521461006%29%28O%29%28747s%29_MAc_.pdf}. The URL is cited as supplied; no PDF has been opened.
\item Literature triage checked 2026-08-21; record updated 2026-09-16. No external source has been visited or independently checked here.
\end{itemize}
\end{document}
